\chapter{Antecedentes}\label{Chap:Antecedentes}
La presente línea de investigación es continuación de dos avances fundamentales; el primero de ellos, presentado por \cite{alcantara} enfocado al análisis de sistemas de levitación magnética,
donde se estudia específicamente el dispositivo ECP 730, para el cual se diseña un control de posición.
Como consecuencia del contacto entre la superficie del imán en forma de disco y la barilla de cristal la cual se desplaza de manera vertical, en dicho dispositivo se presenta el fenómeno de atascamiento-deslizamiento.
Por tanto, se plantea una estrategia de control basada en control conmutado que permite mejorar el desempeño del lazo de control.\\
En el segundo trabajo realizado por \cite{perez}, de igual manera, se manifiesta el efecto de atascamiento-deslizamiento;
sin embargo, en este escenario se estudia un motor de corriente directa(CD) de imanes permanentes, en el que se propone un controlador de doble efector, para controlar la posición del rotor.\\
Es necesario aclarar que en ambos se validó la implemententación de las simulaciones demostrando su factibilidad para realizar físicamente los controladores. 